\section*{Capítulo \ref{ch:b3:endocrinology} --- Endocrinología y homeostasis}

\begin{solution}{exo:b3:endocrinology:1}
\omterm{def:b3:endocrinology:glucose}{Insulina}: péptido, receptor de superficie (una tirosina-cinasa), actúa
en minutos, semivida de unos \qty{5}{min}. Cortisol: esteroide,
\omterm{def:b3:endocrinology:hormone}{receptor nuclear}, horas, semivida de \qty{80}{min} (unido en su mayor
parte a un transportador). Adrenalina: amina, receptores de superficie
acoplados a proteínas G, segundos, semivida de unos \qty{2}{min}.
Tiroxina: aminoácido yodado, \omterm{def:b3:endocrinology:hormone}{receptor nuclear}, días, semivida de
\qty{7}{d} (unida a transportadores). Vasopresina: péptido, receptor de
superficie acoplado a proteínas G (AMPc en el \omterm{def:b3:renal-osmoregulation:nephron}{tubo colector}), minutos,
semivida de unos \qty{15}{min}.
\end{solution}

\begin{solution}{exo:b3:endocrinology:2}
\omterm{def:b3:endocrinology:axes}{Hipotálamo} (TRH) $\to$ \omterm{def:b3:endocrinology:axes}{hipófisis} (TSH) $\to$ tiroides (T$_{4}$), con la
T$_{4}$ inhibiendo los dos pisos superiores. Tiroides destruida:
T$_{4}$ baja, TSH alta. \omterm{def:b3:cancer-biology:hallmarks}{Tumor} secretor de TSH: TSH alta y T$_{4}$ alta
(la retroalimentación no logra suprimir el \omterm{def:b3:cancer-biology:hallmarks}{tumor}). Exceso de pastillas
de tiroxina: T$_{4}$ alta, TSH suprimida. Carencia de yodo: T$_{4}$
baja o en el límite bajo, TSH alta, y la glándula crece bajo ella hasta
formar un bocio.
\end{solution}

\begin{solution}{exo:b3:endocrinology:3}
El encéfalo muere de una glucosa baja en minutos y solo sufre daño por
una alta a lo largo de años; y la evolución encontró la hambruna mucho
más a menudo que los banquetes. De modo que la defensa frente a una
caída es redundante (\omterm{def:b3:endocrinology:glucose}{glucagón}, adrenalina, cortisol, \omterm{def:b3:endocrinology:hormone}{hormona} del
crecimiento) y la defensa frente a una subida es única. Predicción:
demasiada \omterm{def:b3:endocrinology:glucose}{insulina} es letal de forma aguda (coma hipoglucémico), y
demasiado poca es una enfermedad crónica ---que es lo que ve la
clínica.
\end{solution}

\begin{solution}{exo:b3:endocrinology:4}
Un \omterm{def:b3:endocrinology:clock}{zeitgeber} es una señal ambiental que pone en hora un reloj: la luz,
el horario de las comidas, la temperatura, el ejercicio, el calendario
social. El reloj maestro emplea la luz, a través de las \omterm{def:b3:sensory-systems:retina}{células
ganglionares} con melanopsina de la \omterm{def:b3:sensory-systems:retina}{retina}. Tras ocho husos horarios el
reloj se desplaza solo alrededor de una hora al día, de modo que
durante una semana el sueño, el cortisol, la temperatura y el apetito
funcionan con la hora antigua mientras los relojes periféricos,
reajustados por las comidas, derivan a su propio ritmo: el desfase
horario, peor hacia el este, porque adelantar el reloj cuesta más que
retrasarlo.
\end{solution}

\begin{solution}{exo:b3:endocrinology:5}
Respuesta semimáxima con $250$ receptores ocupados: $H = K_{d}\times
250/(\num{20000} - 250) = \qty{1}{nM}\times 0.0127 = \qty{12.7}{pM}$,
ochenta veces por debajo de $K_{d}$. Con \num{2000} receptores: $250/1750
\times \qty{1}{nM} = \qty{143}{pM}$, once veces menos sensible. Una
\omterm{def:b3:endocrinology:glucose}{insulina} alta y prolongada regula a la baja sus receptores, de modo que
un efecto dado exige más \omterm{def:b3:endocrinology:glucose}{insulina}: un componente de la resistencia a la
\omterm{def:b3:endocrinology:glucose}{insulina}, que se alimenta a sí misma.
\end{solution}

\begin{solution}{exo:b3:endocrinology:6}
$T_{4} = 13$: $1.5\,\mathrm{e}^{0.92} = \qty{3.8}{mU/L}$; $11$:
$1.5\,\mathrm{e}^{1.84} = \qty{9.4}{mU/L}$; $9$: $1.5\,\mathrm{e}^{2.76}
= \qty{24}{mU/L}$; $25$: $1.5\,\mathrm{e}^{-4.6} = \qty{0.015}{mU/L}$. A
\qty{11}{pmol/L} la T$_{4}$ sigue dentro de \qtyrange{10}{20}{}
mientras la TSH supera el doble de su límite superior (hipotiroidismo
subclínico); en $13$ la TSH queda justo dentro del intervalo; en $9$
ambas son anormales; en $25$ ambas son anormales del otro modo.
\end{solution}

\begin{solution}{exo:b3:endocrinology:7}
$L = s\beta/(a\gamma) = 0.2\times 0.05/(0.02\times 0.1) = 5$. Exceso
estacionario $g^{*} = (u/a)/(1 + L) = (2/0.02)/6 = \qty{16.7}{mg/dL}$; sin
retroalimentación, $u/a = \qty{100}{mg/dL}$. Exceso de \omterm{def:b3:endocrinology:glucose}{insulina} $i^{*} = \beta g^{*}/
\gamma = 0.05\times 16.7/0.1 = \qty{8.3}{mU/L}$.
\end{solution}

\begin{solution}{exo:b3:endocrinology:8}
Discriminante $(a - \gamma)^{2} - 4s\beta = 0.0064 - 0.04 < 0$:
oscilatorio. $\omega = \sqrt{a\gamma + s\beta - (a+\gamma)^{2}/4} =
\sqrt{0.002 + 0.01 - 0.0036} = \qty{0.092}{min^{-1}}$, período $2\pi/
\omega = \qty{68}{min}$; la amplitud cae en un factor $\mathrm{e}$ en $2/(a +
\gamma) = \qty{16.7}{min}$. Con $s$ a la mitad: $s\beta = 0.005$,
discriminante $0.0064 - 0.02 < 0$, todavía oscilatorio, $\omega = \sqrt{0.007 - 0.0036}
= \qty{0.058}{min^{-1}}$, período \qty{108}{min}; el tiempo de
decaimiento no cambia; $L = 2.5$. Un bucle más débil regresa más
despacio y mantiene un error estacionario mayor.
\end{solution}

\begin{solution}{exo:b3:endocrinology:9}
Durante un año la prednisolona suprimió la CRH y la ACTH; la corteza
suprarrenal, sin estímulo, se atrofió. Al interrumpirla no hay
esteroide exógeno, el \omterm{def:b3:endocrinology:axes}{hipotálamo} y la \omterm{def:b3:endocrinology:axes}{hipófisis} tardan semanas en
reanudarse, y la suprarrenal encogida no podría responder a la ACTH si
llegara. Una infección exige diez veces el cortisol normal; sin
ninguno, los vasos se dilatan y la presión y la glucosa caen: crisis
suprarrenal. Medido: ACTH baja (la lesión es central), cortisol bajo y
una pobre subida del cortisol ante la ACTH inyectada (la glándula se ha
consumido) ---a diferencia de la enfermedad de Addison, donde la ACTH
es alta. De ahí la retirada lenta.
\end{solution}

\begin{solution}{exo:b3:endocrinology:10}
Una variable: $\dot x = f(x)$ con $f' < 0$ tiene un único punto fijo
$x^{*}$; para $x > x^{*}$, $\dot x < 0$ y $x$ decrece monótonamente
hacia $x^{*}$ sin cruzarlo nunca (unicidad de las soluciones), y
simétricamente por debajo: no hay oscilación. Dos variables: la matriz
puede tener autovalores complejos, como en el teorema, cuando el
efector actúa a través de una segunda variable con su propia escala de
tiempo. Una represión instantánea se asentaría en un nivel estable; el
retraso entre la transcripción y la represión nuclear (traducción,
dimerización, fosforilación, importación) deja que la proteína se pase
antes de que el gen se apague y se quede corta antes de que vuelva a
encenderse, y el período es aproximadamente el doble de la suma del
retraso y del tiempo de degradación de la proteína. Una degradación más
rápida de PER acorta el período: una mutación desestabilizadora de PER2
humana da un reloj de unas 20 horas y una familia que se duerme a las
siete de la tarde.
\end{solution}

\begin{solution}{exo:b3:endocrinology:11}
Exponente $k(\text{Ca} - 1.2)$: en $1.10$, $\mathrm{e}^{-2}$, PTH $= 1/
(1 + 0.135) = 0.88$; en $1.15$: $0.73$; en $1.20$: $0.50$; en $1.25$: $1/(1 +
2.72) = 0.27$; en $1.30$: $1/(1 + 7.39) = 0.12$. Pendiente en el \omterm{thm:b3:endocrinology:loop}{punto
de consigna}, $-k/4 = -5$ por mmol/L: un \qty{5}{\%} del máximo por cada
\qty{0.01}{mmol/L}. El calcio ha de mantenerse dentro de un pequeño
porcentaje y su efector actúa sobre un enorme amortiguador (el hueso)
sin pasarse, de modo que una respuesta tan inclinada es segura y hace
falta. Un bucle de glucosa así de inclinado, actuando a través de una
\omterm{def:b3:endocrinology:glucose}{insulina} con su propio tiempo de depuración, tendría una \omterm{thm:b3:endocrinology:loop}{ganancia de
bucle} grande y oscilaría hasta hundirse en la hipoglucemia tras cada
comida.
\end{solution}

\begin{solution}{exo:b3:endocrinology:12}
La glucosilación es irreversible y lenta, de modo que una célula de
edad $\tau$ tiene una fracción $kG\tau$ glucosilada; promediado sobre
células de edades uniformemente repartidas entre \numrange{0}{120}
días, HbA1c $= kG\times\qty{60}{d}$, proporcional a la glucosa
media. Calibración: $5.5\times 60k = 0.05$ da $k =
\qty{1.5e-4}{}$ por (mmol/L)(día); a \qty{10}{mmol/L}: \qty{9.1}{\%}.
Con células que viven 60 días la edad media es de 30 días y la HbA1c se
reduce a la mitad para la misma glucosa: un paciente a \qty{10}{mmol/L}
marcaría un \qty{4.5}{\%}, normal.
\end{solution}

\begin{solution}{pb:b3:endocrinology:1}
\textbf{1.} \qty{90}{mg/dL} $= \qty{0.9}{g/L} = 0.9/180 = \qty{5.0}{mmol/L}$;
$0.9\times 15 = \qty{13.5}{g}$ en el volumen de distribución.
\textbf{2.} $75/15 = \qty{5}{g/L}$: una subida de \qty{500}{mg/dL},
\qty{28}{mmol/L}. El máximo real es mucho menor porque la absorción se
reparte a lo largo de una hora mientras la eliminación retira glucosa a
medida que llega, y el hígado se queda con una buena parte en el primer
paso.
\textbf{3.} $\qty{75}{g}/\qty{60}{min} = \qty{1.25}{g/min}$; sobre
\qty{15}{L}: \qty{8.3}{mg/dL/min}, cuatro veces $u_{0}$.
\textbf{4.} Exceso estacionario $u/a = 8.3/0.02 = \qty{417}{mg/dL}$,
constante de tiempo $1/a = \qty{50}{min}$; al cabo de una hora, $417(1 -
\mathrm{e}^{-1.2}) = \qty{291}{mg/dL}$ por encima del ayuno: unos
\qty{380}{mg/dL}, \qty{21}{mmol/L}.
\textbf{5.} $L = 0.2\times 0.05/(0.02\times 0.1) = 5$; exceso
estacionario $417/6 = \qty{69}{mg/dL}$, seis veces menor que sin
\omterm{def:b3:endocrinology:glucose}{insulina}.
\textbf{6.} El encéfalo quema \qty{120}{g} de glucosa al día, no
almacena nada y poco más puede usar: una glucosa baja da confusión y
coma en minutos. Una glucosa alta glucosila proteínas y, a lo largo de
los años, daña los vasos pequeños de la \omterm{def:b3:sensory-systems:retina}{retina}, el riñón y los nervios,
y las arterias grandes.
\textbf{7.} $\lambda^{2} + (a + \gamma)\lambda + (a\gamma + s\beta) = 0$:
traza $-0.12$, determinante $0.002 + 0.01 = 0.012$, discriminante
$0.0144 - 0.048 = -0.0336$.
\textbf{8.} $\omega = \sqrt{0.0336}/2 = \qty{0.092}{min^{-1}}$; período
\qty{68}{min}; tiempo de decaimiento $2/0.12 = \qty{16.7}{min}$.
\textbf{9.} $\dot g(0) = 100(-0.06 + c\,\omega) = -2$, de modo que $c\,\omega =
0.04$, $c = 0.44$.
\textbf{10.} $g = 0$ cuando $\tan\omega t = -1/c = -2.29$, $\omega t =
\pi - 1.16 = 1.98$, $t = \qty{21.6}{min}$. En el primer mínimo, hacia
\qty{32}{min}: $100\,\mathrm{e}^{-1.94}(\cos 2.97 + 0.44\sin 2.97) =
14.4\times(-0.91) \approx -\qty{13}{mg/dL}$, glucosa \qty{77}{mg/dL}.
\textbf{11.} En $t = 0$, $\dot i = 100\beta > 0$: la \omterm{def:b3:endocrinology:glucose}{insulina} sube
mientras $g > 0$, y alcanza su máximo cuando $\beta g = \gamma i$, lo
que exige que $g$ siga siendo positiva ---de modo que el máximo llega
antes de que la glucosa alcance el valor de ayunas y después de que la
glucosa haya empezado a caer (cae desde $t = 0$). En el modelo, el
máximo queda cerca de los \qty{12}{min}.
\textbf{12.} Con la entrada repartida a lo largo de una hora, el máximo
se alcanza cuando se equilibran la absorción y la eliminación, más bajo
y más tarde (30--60 min); el retorno espera a que termine la absorción,
de ahí las dos horas; el exceso de \omterm{def:b3:endocrinology:glucose}{insulina} es menor con una entrada
gradual, de modo que el descenso es leve.
\textbf{13.} $L = 2.5$. Sano, $g^{*} = (2/0.02)/6 = \qty{16.7}{mg/dL}$;
resistente, $100/3.5 = \qty{28.6}{mg/dL}$: una subida de unos
\qty{12}{mg/dL} (\qty{0.7}{mmol/L}), con la glucosa en ayunas cerca de
\qty{102}{mg/dL}, \qty{5.7}{mmol/L} ---el umbral de la glucemia basal
alterada.
\textbf{14.} $i^{*} = \beta g^{*}/\gamma$: sano, $8.3$; resistente,
\qty{14.3}{mU/L}; razón $1.7$: resistencia a la \omterm{def:b3:endocrinology:glucose}{insulina} compensada,
hiperinsulinemia con una glucosa casi normal.
\textbf{15.} $L = 0.1\times 0.01/0.002 = 0.5$; $g^{*} = 100/1.5 =
\qty{66.7}{mg/dL}$, \qty{50}{mg/dL} por encima del estado sano:
\qty{2.8}{mmol/L}, glucosa en ayunas \qty{7.8}{mmol/L}
(\qty{140}{mg/dL}), por encima del umbral diabético.
\textbf{16.} Discriminante $(a - \gamma)^{2} - 4s\beta = 0.0064 - 0.004 =
0.0024 > 0$: retorno monótono. $\lambda = -0.06 \pm 0.0245$: $-0.0355$
y $-0.0845$ por minuto; la constante de tiempo más lenta es de
\qty{28}{min}, frente a los \qty{16.7}{min} del decaimiento sano: el
bucle que falla regresa más despacio y nunca se pasa por debajo.
\textbf{17.} Glucosa alta con \omterm{def:b3:endocrinology:glucose}{insulina} alta: resistencia con
compensación parcial ---tipo~2; el valor de las dos horas, de
\qty{13}{mmol/L}, supera \qty{11.1}{}, de modo que es diabetes y no una
mera tolerancia alterada.
\textbf{18.} La ausencia de péptido~C significa que no hay \omterm{def:b3:endocrinology:glucose}{insulina}
endógena: tipo~1. Sin \omterm{def:b3:endocrinology:glucose}{insulina}, el músculo y la grasa no pueden captar
glucosa, la grasa se degrada a ácidos grasos y cuerpos cetónicos, la
proteína muscular se cataboliza para la gluconeogénesis y la glucosa se
pierde por la orina con sus calorías: el paciente pasa hambre en medio
de la abundancia.
\textbf{19.} $5\times 9/5.5 = \qty{8.2}{\%}$.
\textbf{20.} $1.5\,\mathrm{e}^{0.46\times 3} = 1.5\times 3.97 =
\qty{6.0}{mU/L}$. T$_{4}$ dentro del intervalo, TSH por encima:
hipotiroidismo subclínico, con la glándula fallando y la \omterm{def:b3:endocrinology:axes}{hipófisis}
compensando.
\textbf{21.} $1.5\,\mathrm{e}^{0.46\times 8} = 1.5\times 39.6 =
\qty{59}{mU/L}$. Una T$_{4}$ baja con una TSH de solo $0.3$ significa
que la \omterm{def:b3:endocrinology:axes}{hipófisis} no responde: la lesión es central (\omterm{def:b3:endocrinology:axes}{hipófisis} o
\omterm{def:b3:endocrinology:axes}{hipotálamo}), no está en la tiroides.
\textbf{22.} Una semivida: el nivel cae al \qty{50}{\%} y los síntomas
se van insinuando a lo largo de semanas. El cortisol, con una semivida
de \qty{80}{min}, desaparece en unas horas tras una dosis olvidada, y
un estrés ese día no encuentra defensa alguna: la semana olvidada de
tiroxina es incómoda, el día olvidado de cortisol puede matar.
\textbf{23.} T$_{4}$ alta (el \omterm{def:b3:adaptive-immunity:antibody}{anticuerpo} impulsa la glándula al margen
de la retroalimentación); TSH suprimida por la T$_{4}$ alta; la
relación log-lineal la lleva hasta lo indetectable ---con $T_{4} = 30$,
$1.5\,\mathrm{e}^{-6.9}
= \qty{0.0015}{mU/L}$--- de modo que una TSH indetectable es el primer
signo.
\textbf{24.} El nivel de una \omterm{def:b3:endocrinology:hormone}{hormona} informa del balance entre su orden
y su glándula: una TSH alta con una T$_{4}$ baja es una glándula que ha
fallado, y una TSH alta con una T$_{4}$ alta, una \omterm{def:b3:endocrinology:axes}{hipófisis}
desbocada; una \omterm{def:b3:endocrinology:glucose}{insulina} alta con una glucosa alta es un cuerpo
resistente, y una \omterm{def:b3:endocrinology:glucose}{insulina} baja con una glucosa alta, un islote
destruido. Leídos a solas, una T$_{4}$ de \qty{7}{pmol/L} o una
\omterm{def:b3:endocrinology:glucose}{insulina} de \qty{30}{mU/L} no dicen dónde está el fallo; leídos junto a
quien los manda, sí.
\textbf{25.} \omterm{thm:b3:endocrinology:loop}{Ganancia de bucle} $L = 5$; período \qty{68}{min} y tiempo
de decaimiento \qty{17}{min}; glucosa en ayunas del bucle que falla,
\qty{7.8}{mmol/L}.
\end{solution}
